% Section 11.2, from lectures/L11/README.md: what you should be able to do afterwards, the
% questions to test yourself, and the next lecture.

\section{Review}
\label{c11:sec:review}

\needspace{6\baselineskip}
\subsection*{What you should be able to do now}
After this chapter, you should be able to:
\begin{itemize}
  \item List four things a subsystem provides that a character device does not.
  \item Say why an \code{ioctl} you invent is a long-term cost, and name the specific portability
    trap.
  \item Implement a \code{gpio_chip} and drive it with \code{libgpiod} tools.
  \item Say what pinctrl is for, and why it is a separate subsystem from GPIO.
  \item Explain why a driver should acquire a clock rather than write a clock gate register.
  \item Say what \code{regmap} replaces, and when it is not worth it.
  \item Register an IIO device with one channel, and read it from sysfs.
  \item Given an unfamiliar device, name the subsystem it belongs to, or say honestly that none
    does.
\end{itemize}

\needspace{6\baselineskip}
\subsection*{Questions to test yourself}
\begin{enumerate}
  \item Your driver exposes a \file{/dev} node and two \code{ioctl}s. Name three things a
    \code{gpiochip} would have given you that you would otherwise write yourself.
  \item Why is \file{/sys/class/gpio} deprecated, and what replaced it?
  \item Two drivers share a clock. One disables it when it suspends. What goes wrong, and which
    subsystem exists to prevent it?
  \item What does \code{regmap} know about your device that your read-modify-write code did not?
  \item An IIO device with a buffer and a trigger delivers samples to userspace. Through what, given
    that you did not implement \code{read}?
  \item You have a device that genuinely fits no subsystem. What do you write, and what do you owe
    the people who will use it?
\end{enumerate}

\needspace{6\baselineskip}
\subsection*{Where to look}
\begin{itemize}
  \item \file{kernel/qa.config} enables \code{CONFIG_GPIO_CDEV} and not \code{CONFIG_GPIO_SYSFS},
    with a comment saying why. That choice is the argument of \secref{c11:app:a5} in one line.
  \item \code{gpioinfo}, \code{gpioget} and \code{gpioset} are not on the target, which carries
    BusyBox and the course's own programs and nothing else. The shipped tester,
    \file{lectures/L11/lab/user/qa_gpio_test.c}, speaks the same ABI they do, and is short enough
    to read in one sitting.
\end{itemize}

\needspace{6\baselineskip}
\subsection*{Looking ahead}
Chapter~\ref{c12:ch} is about:
\begin{itemize}
  \item What all of this costs in latency, and how you would know.
  \item Where latency actually comes from, which is three places and not one.
  \item What \code{PREEMPT_RT} changes about the driver you have just finished.
  \item A measurement that is honest about what it cannot see.
\end{itemize}
